%==============================================================
%  Response to Reviewers' Comments
%  Manuscript ID 88 -- FMLDS2026
%  Page layout and prose follow The Chicago Manual of Style, 17th ed.
%  (manuscript style: 12 pt serif, 1 in margins, double-spaced
%   body text, single-spaced block quotations).
%
%  References use the IEEE numbered style, to match the reference
%  list of the manuscript itself.  Entries that also appear in the
%  manuscript are reproduced from its bibliography verbatim.
%
%  NOTE: All descriptions of changes made to the manuscript are
%  typeset in RED, via the \begin{changesblock} ... \end{changesblock}
%  environment defined below.
%==============================================================

\documentclass[12pt,letterpaper]{article}

% ---------- Page layout (Chicago: 1-inch margins all round) ----------
\usepackage[margin=1in]{geometry}

% ---------- Fonts (Chicago: a readable serif, Times New Roman) -------
\usepackage{mathptmx}
\usepackage[T1]{fontenc}
\usepackage[utf8]{inputenc}

% ---------- Spacing (Chicago manuscript style: double-spaced) --------
\usepackage{setspace}
\doublespacing
% To produce a shorter document, replace \doublespacing above with
% \onehalfspacing.  Nothing else needs to change.

% ---------- Utility packages ----------------------------------------
\usepackage{amsmath}
\usepackage{amssymb}
\usepackage{xcolor}
\usepackage{enumitem}
\usepackage{fancyhdr}
\usepackage[hidelinks]{hyperref}

% ---------- The colour used to mark changes --------------------------
% Defined once here so it can be changed globally in a single place.
\definecolor{changered}{rgb}{0.80,0.00,0.00}

% ---------- Running head (Chicago: author surname + page no.) --------
\pagestyle{fancy}
\fancyhf{}
\renewcommand{\headrulewidth}{0pt}
\rhead{\small Meruva et al.\ \thepage}

% ---------- Sectional headings (Chicago levels 1 and 2) --------------
\usepackage{titlesec}
\titleformat{\section}{\normalfont\large\bfseries\centering}{}{0pt}{}
\titleformat{\subsection}{\normalfont\normalsize\bfseries}{}{0pt}{}
\titlespacing*{\section}{0pt}{1.5\baselineskip}{0.5\baselineskip}
\titlespacing*{\subsection}{0pt}{1\baselineskip}{0.25\baselineskip}

% ---------- Environment for verbatim reviewer comments ---------------
% Chicago: block quotations are single-spaced, indented, and set
% without quotation marks.
\newenvironment{reviewercomment}
  {\begin{singlespace}%
   \begin{quote}\itshape\noindent}
  {\end{quote}%
   \end{singlespace}\vspace{-0.25\baselineskip}}

% ---------- Environment for changes made (typeset in red) ------------
% An environment (rather than a command) is used so that the colour
% change is confined to its own group and cannot bleed into the text
% that follows.
\newenvironment{changesblock}
  {\par\vspace{0.3\baselineskip}%
   \color{changered}\noindent
   \textbf{Changes to the manuscript.}\enspace\ignorespaces}
  {\par\vspace{0.3\baselineskip}}

% ---------- Label for each numbered reviewer point -------------------
\newcommand{\point}[1]{%
  \par\normalcolor\vspace{0.4\baselineskip}%
  \noindent\textbf{#1}\par\nopagebreak}
\newcommand{\response}{\noindent{\color{changered}\textbf{Response.}}\enspace}

% ---------- IEEE-style numbered reference list -----------------------
% Matches the numbered "[n]" bibliography style of the manuscript.
\newenvironment{ieeerefs}
  {\begin{singlespace}%
   \begin{list}{[\arabic{enumi}]}{%
     \usecounter{enumi}%
     \setlength{\labelwidth}{2em}%
     \setlength{\labelsep}{0.5em}%
     \setlength{\leftmargin}{2.5em}%
     \setlength{\itemindent}{0pt}%
     \setlength{\listparindent}{0pt}%
     \setlength{\topsep}{0.5\baselineskip}%
     \setlength{\itemsep}{0.5\baselineskip}%
     \setlength{\parsep}{0pt}}}
  {\end{list}\end{singlespace}}

\begin{document}

%==============================================================
%  Heading
%==============================================================
\begin{singlespace}
\begin{center}
{\large\bfseries Response to Reviewers' Comments}\\[0.5\baselineskip]
{\normalsize\bfseries Manuscript ID 88}\\[0.25\baselineskip]
{\normalsize ``Uncertainty-Aware Multimodal Deep Learning Framework for\\
Stress Detection Using Audio-Visual Signals''}\\[0.5\baselineskip]
{\normalsize Track: FMLDS2026}\\[0.75\baselineskip]
{\normalsize Authors - Meruva Kodanda Suraj, Mamta Nag, and Mallikharjuna Rao K.}\\[0.25\baselineskip]
{\normalsize International Institute of Information Technology, Naya Raipur}\\[0.25\baselineskip]
{\normalsize Chhattisgarh, India}
\end{center}
\end{singlespace}

\vspace{0.5\baselineskip}

\begin{singlespace}
\noindent\fbox{\parbox{\dimexpr\linewidth-2\fboxsep-2\fboxrule\relax}{%
\small\textbf{Note on presentation.}\enspace Reviewers' comments are
reproduced verbatim as indented italic block quotations. {\color{changered}%
Section headings and all descriptions of changes made to the manuscript
are set in red}, so that the replies and the specific revisions can be
located at a glance. References follow the numbered IEEE style used in
the manuscript; see the note preceding the reference list.}}
\end{singlespace}

\vspace{0.5\baselineskip}

%==============================================================
%  Cover statement
%==============================================================
\section*{To the Program Chairs and Reviewers}
\addcontentsline{toc}{section}{To the Program Chairs and Reviewers}

We are grateful to both reviewers for the care and specificity of their
reports. The comments identified genuine defects in the submitted
manuscript---an unsubstantiated architectural claim, a numerical
inconsistency between a table and a figure, several incomplete captions,
and the omission of the experimental detail required for
reproducibility---and correcting them has materially improved the paper.
We have revised the manuscript throughout and respond to every point
below.

Reviewer 3's comments are addressed in the first section and Reviewer 5's
in the second. Each comment is reproduced verbatim as a block quotation,
followed by our response and a statement of the specific changes made to
the manuscript. Section, table, and figure numbers refer to the revised
manuscript unless stated otherwise. Where a reviewer's request could not
be satisfied within the scope of a conference revision, we say so plainly
rather than claim otherwise, and we describe the interim evidence we have
supplied in its place.

Before turning to the individual points, we summarize the four
substantive changes that the revision makes.

\begin{singlespace}
\color{changered}
\begin{enumerate}[leftmargin=*,itemsep=0.4\baselineskip]

\item \textbf{The Action Unit gating claim has been withdrawn in full.}
The component was not implemented in the evaluated model, and the claim
should not have appeared in the submission. Every mention has been
removed from the abstract, the introduction, the architecture
description, and the contribution list, and Figure~1 has been redrawn so
that the diagram matches the architecture actually evaluated.

\item \textbf{The contradiction between the confusion matrix and the
per-class report has been resolved.} The two class rows in the per-class
table had been transposed. The confusion matrix was correct; the table
was not. The table now reports the values that follow from the matrix.

\item \textbf{A complete experimental configuration is now reported.}
A new Table~I gives the dataset size, the split sizes and protocol, the
optimizer, learning rate, batch size, schedule, epoch budget, dropout
rate, number of Monte Carlo passes, temperature, precision, and random
seed.

\item \textbf{The accuracy of the proposed model is now reported with an
interval, and the limits of the evidence are stated openly.} We report
Wilson confidence intervals for the test accuracy, state that they
overlap those of the CNN-LSTM baseline, disclose that the split is not
actor-disjoint, and characterize the uncertainty AUROC of 0.6906 as a
weak-to-moderate discriminator rather than a strength of the method.

\end{enumerate}
\end{singlespace}

%==============================================================
%  REVIEWER 3
%==============================================================
\section*{Response to Reviewer 3}
\addcontentsline{toc}{section}{Response to Reviewer 3}

We thank Reviewer 3 for a detailed and exacting report. The reviewer's
overall assessment was Strong Reject, and we accept that the submitted
version did not meet the standard of technical correctness the report
demands. We have treated each of the ten points as a defect to be
corrected rather than a matter for argument.

%--------------------------------------------------------------
\point{Comment 3.1 (Major). Action Unit gating is missing from the paper.}

\begin{reviewercomment}
In the abstract, AU gating is listed as a contribution criterion.
However, Section III does not give any description, formula, or
architectural detail. Please either provide a complete description,
mathematical formulation, or remove it completely from all claims. As
mentioned in the abstract, the contribution can't be verified.
\end{reviewercomment}

\response The reviewer is correct, and we accept the criticism without
qualification. The Action Unit gating mechanism was not implemented in
the model that produced the reported results. Its appearance in the
abstract and introduction of the submitted manuscript was an error
carried over from an earlier design of the architecture, and it described
a component that no reported number depends on. Because the claim cannot
be verified---as the reviewer observes---and because we do not have
Facial Action Coding System~[1] annotations or an AU
detector integrated into the present pipeline, we have taken the
reviewer's second option and removed the claim entirely rather than
retrofit a formulation after the fact.

\begin{changesblock}
Every reference to Action Unit gating has been deleted from the
abstract, the introduction, the contribution list, and Section~III.
The architecture is now described consistently as a hybrid of a 3D
convolutional network, EfficientNet-B0~[2], an audio
convolutional branch, and cross-modal attention fusion---the four
components that are in fact implemented and evaluated. Figure~1 has been
redrawn accordingly; see also our response to Comment~5.4. The
contribution list in Section~I now names only cross-modal attention
fusion, Monte Carlo Dropout and temperature scaling, the calibration and
uncertainty evaluation, and selective prediction.
\end{changesblock}

%--------------------------------------------------------------
\point{Comment 3.2 (Major). The proposed model underperforms the baseline model.}

\begin{reviewercomment}
UA-HybridNet (Acc: 0.813, AUC: 0.878) is outperformed by CNN-LSTM (Acc:
0.885, AUC: 0.971) on every metric. The authors argue that uncertainty
justifies this trade-off, which requires proof. Please show: (a) that the
CNN-LSTM baseline is poorly calibrated (report its ECE/MCE), and (b) that
applying temperature scaling to the CNN-LSTM does not close the
calibration gap while retaining its higher accuracy. Without this, it is
unclear why UA-HybridNet is preferable.
\end{reviewercomment}

\response This is the most consequential comment in either report, and we
answer it directly. The reviewer has identified that our justification
for the accuracy trade-off was asserted rather than demonstrated, and
that the decisive experiment---calibrating the stronger baseline and
seeing whether it then dominates on both axes---was not run. We agree
that this experiment is the right test of our claim, and we agree that
without it the preference for UA-HybridNet is not established. We are
unable to report the paired ECE and MCE for a temperature-scaled CNN-LSTM
within this revision cycle, because doing so properly requires refitting
the baseline's temperature on a validation split under the same
actor-disjoint protocol discussed in Comment~3.3, and reporting a
hurried result would repeat the error the reviewer is objecting to.

We have therefore withdrawn the unsupported claim rather than defend it,
and have made three changes that we hope the reviewer will accept as an
honest interim position. First, the manuscript no longer asserts that the
uncertainty machinery justifies the accuracy gap; the claim is now
confined to what the reported evidence supports, namely that temperature
scaling reduces UA-HybridNet's own expected calibration error from 0.1768
to 0.0710 and its maximum calibration error from 0.7651 to 0.2531
(Table~IV), and that selective prediction converts uncertainty estimates
into a measurable reduction in risk (Table~VI). Second, we now report
that the accuracy difference is not statistically established on a test
set of this size: the 95 percent Wilson interval~[3] for
UA-HybridNet's accuracy is $[0.743, 0.868]$ against approximately
$[0.826, 0.928]$ for the baseline, and these intervals overlap. Third, we
state in Section~IV-D that fitting temperature scaling to the baseline's
own logits and reporting its ECE and MCE alongside ours is the first of
our planned extensions, and we identify it by name as the experiment
needed to settle the comparison.

We recognize that this is a weaker answer than the reviewer asked for. We
judged it better to mark the claim as unproven and name the experiment
that would prove it than to present an underpowered result as
confirmation.

\begin{changesblock}
The discussion following Table~II has been rewritten to remove
the claim that reliability compensates for the accuracy deficit.
Section~IV-D now opens with the baseline calibration experiment as the
first planned extension, stating that we will fit temperature scaling to
the CNN-LSTM's logits and report its ECE and MCE alongside those in
Table~IV ``rather than leaving the comparison one-sided.'' The Wilson
intervals for both models are reported in the same subsection, with the
baseline's success count inferred from its reported accuracy and recorded
in a footnote.
\end{changesblock}

%--------------------------------------------------------------
\point{Comment 3.3 (Major). Artificial actors in RAVDESS and lack of stress data.}

\begin{reviewercomment}
RAVDESS is a set of emotional data recorded by human actors, not actual
stress data. Further, there are 24 actors who performed emotional
expressions for the recordings, not acted on behalf of emotional data
recording. Make sure your data is actor-disjoint such that no actor
appears in both the training and test sets. Else, the model may learn
speaker identity rather than stress.
\end{reviewercomment}

\response We accept both halves of this comment. On the first, RAVDESS
consists of portrayed emotional expressions recorded by professional
actors under controlled conditions~[4]; it is
not a corpus of spontaneous physiological or psychological stress, and
our mapping of anger, fear, and sadness to a ``stress'' class is a
proxy, not a clinical label. The manuscript now says this in those terms.

On the second, the reviewer's diagnosis is correct and we did not
implement an actor-disjoint split. Our 70/15/15 partition was performed
at the clip level by stratified random sampling, which means that clips
from the same actor can appear in more than one split. The consequence
the reviewer names---that the model may key on speaker identity rather
than on stress-related expression---is a real threat to the validity of
our accuracy figures, and we cannot rule it out from the experiments
reported here. Re-running every model under a leave-actors-out protocol
would change all reported numbers and is beyond what we can complete for
this revision.

Rather than leave this implicit, we have made the leakage risk explicit
in three places in the manuscript so that no reader can mistake our
results for subject-independent performance. We consider a leave-actors-out
re-evaluation to be the second priority for the extended version of this
work, after the baseline calibration experiment of Comment~3.2.

\begin{changesblock}
Section~IV-A now states that ``Partitioning was performed at the
clip level by stratified random sampling, so clips from the same actor
may span splits; the evaluation is thus not strictly subject-independent.''
The Limitations subsection repeats the point and adds that a
leave-actors-out protocol is planned. Section~IV-D names the
re-evaluation as a planned extension whose purpose is ``to rule out
speaker-identity leakage and confirm the reported accuracy holds for
unseen speakers.'' The Limitations subsection also now states that
RAVDESS ``includes performed emotional expressions rather than
spontaneous real-world stress signals,'' which may affect applicability
to real-world settings.
\end{changesblock}

%--------------------------------------------------------------
\point{Comment 3.4 (Major). Too small a data size.}

\begin{reviewercomment}
The test set has only 150 samples (75 per class). At this scale, the
confidence interval will be very large, which remains unreported in the
paper. Minor variance in the test set might shift accuracy by several
percentage points.
\end{reviewercomment}

\response We agree, and the reviewer's estimate of the effect is
accurate. On 150 test samples the sampling error is large enough that
differences of several percentage points carry little evidential weight,
and the submitted manuscript reported point estimates as though they were
precise. We have corrected this by reporting interval estimates.

Using the Wilson score interval~[3], which is preferred over
the normal approximation at this sample size, UA-HybridNet's accuracy of
122/150 = 0.8133 has a 95 percent interval of $[0.743, 0.868]$, a width
of roughly 12.5 percentage points. The corresponding interval for the
CNN-LSTM baseline is approximately $[0.826, 0.928]$. As noted in our
response to Comment~3.2, these intervals overlap, so the apparent
ordering of the two models is not resolved by this test set. We now say
so in the manuscript instead of comparing point estimates as if the gap
were established.

\begin{changesblock}
Section~IV-D reports both Wilson intervals and states that the
modest test set motivates expanding the evaluation set and running a
paired McNemar test~[5] on matched predictions ``to confirm
the accuracy gap against the baseline is real.'' The Limitations
subsection notes that the dataset is modest relative to comparable
deep learning corpora.
\end{changesblock}

%--------------------------------------------------------------
\point{Comment 3.5 (Major). Missing experimental details.}

\begin{reviewercomment}
Please report the total dataset size, train size, validation size, test
size, actor overlap policy, MC Dropout rate, number of training epochs,
optimizer, learning rate, and batch size. The submitted paper is not
reproducible.
\end{reviewercomment}

\response We accept this without reservation; the submitted version did
not contain enough information to reproduce the results. Every item on
the reviewer's list is now reported.

\begin{changesblock}
A new Section~IV-A (``Experimental Setup'') and an accompanying
Table~I have been added. Table~I reports the total sample count (1000,
balanced 500/500), the train, validation, and test sizes (700/150/150)
under a 70/15/15 stratified split with seed 42, the optimizer (AdamW with
weight decay $10^{-5}$), the learning rate ($10^{-4}$), the batch size (4,
effective 16 through gradient accumulation), the schedule (cosine
annealing), the epoch budget (maximum 100 with early-stopping patience
20), the dropout rate (0.3), the number of Monte Carlo Dropout passes
(20), the number of test-time augmentation copies (5), the fitted
temperature ($T = 3.0855$), and the numerical precision (16-bit mixed).
The actor overlap policy---clip-level stratified sampling, not
actor-disjoint---is stated in the running text of Section~IV-A, as
described under Comment~3.3. Section~III-B additionally now specifies the
segmentation procedure that produces the 1000 samples: fixed-length
two-second clips extracted with a 0.5-second sliding-window stride.
\end{changesblock}

%--------------------------------------------------------------
\point{Comment 3.6 (Minor). Contradictory Table II caption.}

\begin{reviewercomment}
Table II caption states, ``the deterministic model gives calibrated
confidence estimates'' --- this directly contradicts the paper's
argument. Please correct it to ``the deterministic model does NOT give
calibrated confidence estimates.'' Else, the claim contradicts the
paper's claim.
\end{reviewercomment}

\response The reviewer is right; the sentence stated the opposite of what
we meant and of what the Reliability column of the same table records.
The negation was dropped in drafting. We have corrected it in the sense
the reviewer indicates.

\begin{changesblock}
The passage now reads: ``the deterministic Hybrid model attains
higher raw accuracy but gives no calibrated confidence estimates, which
are necessary in safety-critical applications for reliable deployment.''
This is consistent with the Reliability column of Table~II, which records
``No'' for all three deterministic architectures and ``Yes'' only for the
uncertainty-aware model.
\end{changesblock}

%--------------------------------------------------------------
\point{Comment 3.7 (Minor). Incomplete Figure 2 caption.}

\begin{reviewercomment}
Figure 2 reads ``Confusion Matrix for .'' --- the model name is missing.
Please complete this caption.
\end{reviewercomment}

\response Corrected. The caption had been left with an unfilled model
name.

\begin{changesblock}
The caption of Figure~2 now reads ``Confusion Matrix for
UA-HybridNet.'' We have also audited every other caption in the
manuscript for the same class of omission; the per-class table
(Table~III) is now captioned ``Per-Class Classification Report for
UA-HybridNet,'' and the model is named consistently in all captions,
table headings, and cross-references.
\end{changesblock}

%--------------------------------------------------------------
\point{Comment 3.8 (Minor). Uncertainty AUROC barely above chance.}

\begin{reviewercomment}
The Uncertainty AUROC of 0.6906 is barely above chance (0.5). It suggests
uncertainty estimation is a weak indicator of misclassification. Please
discuss the result honestly.
\end{reviewercomment}

\response We agree with the reviewer's reading of the number and have
adopted it in the manuscript. An AUROC of 0.6906 for separating correct
from incorrect predictions by predictive uncertainty is a
weak-to-moderate discriminator, and it does not support a strong claim
that Monte Carlo Dropout reliably identifies the model's errors. The
submitted version reported the value without comment, which invited a
more favorable reading than it deserves.

We note, without offering it as a defense, that the selective prediction
results in Table~VI are consistent with a weak but non-zero signal:
rejecting the 20 percent most uncertain samples raises accuracy from
0.8133 to 0.8333 and AUC from 0.8782 to 0.8999, and raises stress recall
from 0.8667 to 0.8871. These are modest gains, of the size one would
expect from a modestly informative uncertainty estimate, and the
manuscript now presents them as such. The natural inference is that the
uncertainty signal itself needs to be improved rather than that the
selective mechanism has been vindicated, and we say this explicitly.

\begin{changesblock}
Section~IV-D now states that ``the current Uncertainty AUROC of
0.6906 is only a weak-to-moderate discriminator of misclassification''
and identifies strengthening the uncertainty estimate as a planned
extension, specifically by exploring deep ensembles~[6]
and evidential deep learning~[7] as alternatives to Monte
Carlo Dropout~[8]. The claims made for the
uncertainty results elsewhere in Section~IV have been moderated to match.
\end{changesblock}

%--------------------------------------------------------------
\point{Comment 3.9 (Minor). Standard equations.}

\begin{reviewercomment}
Equations 1--5 are standard from previous work. Please consider replacing
those with your own mathematical formulation.
\end{reviewercomment}

\response We have partly followed this suggestion and should be candid
about the extent. The underlying quantities---Monte Carlo Dropout
statistics, predictive entropy and mutual information, the
temperature-scaling objective, and the expected calibration error---are
established results~[8]--[10], and we do not think it would be honest to present
re-derivations of them as our own contribution. What we have done instead
is rewrite the formal section so that the notation is our own and
specific to this architecture, and so that each equation states something
the reader needs in order to follow our system rather than reciting
background.

Concretely, the fusion stage is now given a compact operator definition,
$\mathrm{CA}(Q,K) = \mathrm{LN}\!\left(Q + \mathrm{MHA}(Q,K,K)\right)$,
which is applied in both directions to yield an audio-aware visual vector
and a visual-aware audio vector, with the modality shapes
($F_v \in \mathbb{R}^{L_v \times d}$, $F_a \in \mathbb{R}^{L_a \times d}$,
and the pooled vectors $f_s, f_v, f_a \in \mathbb{R}^{d}$) defined
explicitly. The temperature-scaling objective is stated as the negative
log-likelihood minimization actually solved on our validation set, with
our fitted value $T^{*} = 3.0855$ reported. The selective-prediction rule
and the calibration error are combined into a single equation block that
defines the abstention threshold $\eta_\alpha$ at coverage $1 - \alpha$
alongside the ECE, because in our framework these two are used together.
We believe this makes Section~III specific to the proposed system without
claiming credit for standard results.

\begin{changesblock}
Section~III-D through III-F have been rewritten with consistent
notation defined at the point of first use. The cross-modal attention
operator, the shapes of all intermediate representations, the fitted
temperature, and the combined selective-prediction and calibration block
are new. Attribution to prior work is given inline for each standard
quantity.
\end{changesblock}

%--------------------------------------------------------------
\point{Comment 3.10 (Minor). Grammatical errors.}

\begin{reviewercomment}
Several grammatical errors reduce readability. The paper would benefit
from careful proofreading.
\end{reviewercomment}

\response Accepted. The manuscript has been proofread in full. Reviewer 5
identified one instance specifically (``competitive perform''), which is
addressed under Comment~5.5.

\begin{changesblock}
The abstract has been rewritten, and the introduction,
Section~III, and the results discussion have been substantially edited
for grammar, article usage, subject-verb agreement, and consistency of
terminology. The model name has been standardized as ``UA-HybridNet''
throughout, having previously alternated between several informal
designations.
\end{changesblock}

%==============================================================
%  REVIEWER 5
%==============================================================
\section*{Response to Reviewer 5}
\addcontentsline{toc}{section}{Response to Reviewer 5}

We thank Reviewer 5 for a concise report that isolated a serious
numerical error we had not detected. Comments 5.1 and 5.2 concern the
same defect and are answered together.

%--------------------------------------------------------------
\point{Comments 5.1 and 5.2. Contradiction between the confusion matrix and the per-class report.}

\begin{reviewercomment}
There is a major mathematical contradiction between the confusion matrix
in Figure 2 and the per-class classification report in Table III.

The calculated recall for the stress class from the confusion matrix is
approximately 86.7\%, but Table III lists it as 76.0\%. The calculated
recall for the non-stress class from the confusion matrix is 76.0\%, but
Table III lists it as 87.0\%
\end{reviewercomment}

\response The reviewer is correct in every particular, including the
arithmetic. We are grateful for the catch. The cause was a transposition
of the two class rows when the classification report was transcribed into
Table~III: the values were individually right but assigned to the wrong
classes, which inverted the interpretation of the model's behavior. The
confusion matrix was correct; the table was not.

The confusion matrix contains 57 true non-stress, 18 non-stress
misclassified as stress, 10 stress misclassified as non-stress, and 65
true stress, over 75 samples per class. This yields stress recall
$65/75 = 0.8667$ and non-stress recall $57/75 = 0.7600$, exactly as the
reviewer computed, together with stress precision $65/83 = 0.7831$ and
non-stress precision $57/67 = 0.8507$. Table~III now reports these
values. All four numbers are mutually consistent with the matrix, with
the overall accuracy of $122/150 = 0.8133$ in Table~II, and with the
0 percent rejection row of the selective-prediction table (Table~VI).

The correction also changes what the paper says about the model, and we
have updated the narrative accordingly. The submitted version, following
the transposed table, described the model as more precise than sensitive
on the stress class. The correct reading is the opposite and is more
favorable to our stated motivation: the model is more sensitive than
precise on stress, with recall 0.8667 against precision 0.7831, which is
the error profile one wants in a monitoring application where a missed
stress case is costlier than a false alarm. We now make this point using
the corrected figures.

\begin{changesblock}
Table~III has been corrected: the Stress row now reads precision
0.7831, recall 0.8667, F1 0.8228; the Non-Stress row reads precision
0.8507, recall 0.7600, F1 0.8028. The macro and weighted averages
(0.8169, 0.8133, 0.8128) are recomputed and, on a balanced test set,
coincide. The accompanying text now states that ``the stress class
achieves higher recall (0.8667), indicating strong sensitivity in
detecting stress, while the non-stress class achieves higher precision
(0.8507).'' The confusion matrix discussion states the correct totals
(122 correct, 28 misclassified). We have verified every derived quantity
in Section~IV against the confusion matrix counts.
\end{changesblock}

%--------------------------------------------------------------
\point{Comment 5.3. Undefined ``Pipeline D'' and incomplete caption.}

\begin{reviewercomment}
Figure 2 contains the header ``Pipeline D Confusion Matrix,'' but the term
``Pipeline D'' is never introduced or explained in the manuscript.Also the
text caption is also incomplete
\end{reviewercomment}

\response Corrected. ``Pipeline D'' was internal shorthand from our
experiment tracking---the fourth configuration in our ablation
sequence---and it should never have reached the manuscript. It denotes
the same model the paper calls UA-HybridNet.

\begin{changesblock}
Figure~2 has been regenerated without the embedded plot title,
so the figure no longer carries the ``Pipeline D'' header. The caption is
now complete and reads ``Confusion Matrix for UA-HybridNet.'' We have
searched the manuscript for other internal designations of this kind;
Table~III, which was similarly affected, is now captioned ``Per-Class
Classification Report for UA-HybridNet.'' The name ``UA-HybridNet'' is
used consistently in all text, captions, and figures. This point overlaps
with Reviewer 3's Comment~3.7, and the two were corrected together.
\end{changesblock}

%--------------------------------------------------------------
\point{Comment 5.4. Action Unit gating absent from Figure 1.}

\begin{reviewercomment}
The abstract and introduction highlight an Action Unit gating system as a
core part of the hybrid architecture, but this component is completely
missing from the system workflow diagram in Figure 1.
\end{reviewercomment}

\response The reviewer has identified the same inconsistency as
Reviewer 3's Comment~3.1, and from the opposite direction: the diagram
was right and the text was wrong. Action Unit gating was absent from
Figure~1 because it was absent from the implemented system. We have
resolved the inconsistency by removing the claim from the text rather
than by adding a component to the diagram.

\begin{changesblock}
All references to Action Unit gating have been removed from the
abstract, introduction, contribution list, and Section~III. Figure~1 has
been redrawn and now shows exactly the evaluated pipeline: the RAVDESS
source, the parallel video and audio preprocessing paths (frame
extraction, face detection, resize and normalize; audio signal, Mel
spectrogram, normalization), the three feature branches (3D-CNN,
EfficientNet-B0, audio CNN), cross-modal attention, multimodal feature
fusion, the fully connected layer, and the binary stress/non-stress
output. Figure~1, Section~III, and the abstract now describe the same
architecture, and no component appears in one but not the others.
\end{changesblock}

%--------------------------------------------------------------
\point{Comment 5.5. Typographical and grammatical errors in the abstract.}

\begin{reviewercomment}
The abstract contains noticeable typographical and grammatical errors,
such as using the phrase ``competitive perform'' instead of ``competitive
performance''.
\end{reviewercomment}

\response Corrected, and the abstract has been rewritten in full rather
than patched. As noted under Comment~3.10, the remainder of the
manuscript has been proofread as well.

\begin{changesblock}
The specific phrase the reviewer cites now reads ``a
competitive performance of 81.3\% accuracy and an AUC of 0.878.'' The
rewritten abstract also removes the withdrawn Action Unit gating claim
(Comment~5.4) and states the contribution in terms of the components
actually evaluated.
\end{changesblock}

%==============================================================
%  Closing
%==============================================================
\section*{Concluding Remarks}
\addcontentsline{toc}{section}{Concluding Remarks}

We have accepted every correction the reviewers identified. The
unverifiable architectural claim is withdrawn, the numerical
contradiction between Table~III and Figure~2 is resolved in favor of the
confusion matrix, the incomplete captions and the internal ``Pipeline D''
designation are removed, and a full experimental configuration is now
reported in Table~I.

Two of Reviewer 3's requests---the calibration comparison against a
temperature-scaled CNN-LSTM baseline (Comment~3.2) and the
actor-disjoint re-evaluation (Comment~3.3)---are not satisfied by new
experiments in this revision. We have not papered over this. In both
cases we have removed the claims that depended on the missing evidence,
stated the limitation in the body of the paper rather than only in a
closing paragraph, and named the specific experiment required to settle
the question. On the related question of statistical power
(Comment~3.4), we now report Wilson intervals that overlap between the
proposed model and the baseline, and we say in the manuscript that the
accuracy ordering is therefore not established by this test set.

We hope the reviewers find that the revised manuscript is internally
consistent, reproducible from the reported configuration, and candid
about the boundary between what the experiments demonstrate and what
remains to be shown. We are grateful for the reviews, which were
unusually specific and correspondingly useful.

%==============================================================
%  References (Chicago author-date)
%==============================================================
\section*{References}
\addcontentsline{toc}{section}{References}

\begin{singlespace}
\noindent\small Reference numbers below are local to this letter and are
assigned in order of first citation; they do not correspond to the
numbering of the manuscript's own bibliography. Entries~[2], [4], [6],
and [8]--[10] also appear in the manuscript and are reproduced from it
verbatim. Entries~[1], [3], [5], and [7] are cited only in this letter.
\end{singlespace}

\vspace{0.5\baselineskip}

\begin{ieeerefs}

\item
P. Ekman and W. V. Friesen,
\emph{Facial Action Coding System: A Technique for the Measurement of Facial Movement}.
Palo Alto, CA: Consulting Psychologists Press, 1978.

\item
M. Tan and Q. V. Le,
``EfficientNet: Rethinking model scaling for convolutional neural networks,''
Proc. International Conference on Machine Learning (ICML), 2019.

\item
E. B. Wilson,
``Probable inference, the law of succession, and statistical inference,''
Journal of the American Statistical Association, vol. 22, no. 158, pp. 209--212, 1927.

\item
S. R. Livingstone and F. A. Russo,
``The Ryerson Audio-Visual Database of Emotional Speech and Song (RAVDESS): A dynamic multimodal database of emotional expressions,''
PLoS ONE, vol. 13, no. 5, 2018.

\item
Q. McNemar,
``Note on the sampling error of the difference between correlated proportions or percentages,''
Psychometrika, vol. 12, no. 2, pp. 153--157, 1947.

\item
B. Lakshminarayanan, A. Pritzel, and C. Blundell,
``Simple and scalable predictive uncertainty estimation using deep ensembles,''
Advances in Neural Information Processing Systems (NeurIPS), 2017.

\item
M. Sensoy, L. Kaplan, and M. Kandemir,
``Evidential deep learning to quantify classification uncertainty,''
Advances in Neural Information Processing Systems (NeurIPS), 2018.

\item
Y. Gal and Z. Ghahramani,
``Dropout as a Bayesian approximation: Representing model uncertainty in deep learning,''
Proc. International Conference on Machine Learning (ICML), 2016.

\item
C. Guo, G. Pleiss, Y. Sun, and K. Q. Weinberger,
``On calibration of modern neural networks,''
Proc. International Conference on Machine Learning (ICML), 2017.

\item
A. Kendall and Y. Gal,
``What uncertainties do we need in Bayesian deep learning for computer vision?''
Advances in Neural Information Processing Systems (NeurIPS), 2017.

\end{ieeerefs}

\end{document}
